%Version 3.1 December 2024
% See section 11 of the User Manual for version history
%
%%%%%%%%%%%%%%%%%%%%%%%%%%%%%%%%%%%%%%%%%%%%%%%%%%%%%%%%%%%%%%%%%%%%%%
%%                                                                 %%
%% Please do not use \input{...} to include other tex files.       %%
%% Submit your LaTeX manuscript as one .tex document.              %%
%%                                                                 %%
%% All additional figures and files should be attached             %%
%% separately and not embedded in the \TeX\ document itself.       %%
%%                                                                 %%
%%%%%%%%%%%%%%%%%%%%%%%%%%%%%%%%%%%%%%%%%%%%%%%%%%%%%%%%%%%%%%%%%%%%%

%%\documentclass[referee,sn-basic]{sn-jnl}% referee option is meant for double line spacing

%%=======================================================%%
%% to print line numbers in the margin use lineno option %%
%%=======================================================%%

%%\documentclass[lineno,pdflatex,sn-basic]{sn-jnl}% Basic Springer Nature Reference Style/Chemistry Reference Style

%%=========================================================================================%%
%% the documentclass is set to pdflatex as default. You can delete it if not appropriate.  %%
%%=========================================================================================%%

%%\documentclass[sn-basic]{sn-jnl}% Basic Springer Nature Reference Style/Chemistry Reference Style

%%Note: the following reference styles support Namedate and Numbered referencing. By default the style follows the most common style. To switch between the options you can add or remove “Numbered” in the optional parenthesis. 
%%The option is available for: sn-basic.bst, sn-chicago.bst%  
 
%%\documentclass[pdflatex,sn-nature]{sn-jnl}% Style for submissions to Nature Portfolio journals
%%\documentclass[pdflatex,sn-basic]{sn-jnl}% Basic Springer Nature Reference Style/Chemistry Reference Style
\documentclass[pdflatex,sn-mathphys-num]{sn-jnl}% Math and Physical Sciences Numbered Reference Style
%%\documentclass[pdflatex,sn-mathphys-ay]{sn-jnl}% Math and Physical Sciences Author Year Reference Style
%%\documentclass[pdflatex,sn-aps]{sn-jnl}% American Physical Society (APS) Reference Style
%%\documentclass[pdflatex,sn-vancouver-num]{sn-jnl}% Vancouver Numbered Reference Style
%%\documentclass[pdflatex,sn-vancouver-ay]{sn-jnl}% Vancouver Author Year Reference Style
%%\documentclass[pdflatex,sn-apa]{sn-jnl}% APA Reference Style
%%\documentclass[pdflatex,sn-chicago]{sn-jnl}% Chicago-based Humanities Reference Style

%%%% Standard Packages
%%<additional latex packages if required can be included here>

\usepackage{graphicx}%
\usepackage{multirow}%
\usepackage{amsmath,amssymb,amsfonts}%
\usepackage{amsthm}%
\usepackage{mathrsfs}%
\usepackage[title]{appendix}%
\usepackage{xcolor}%
\usepackage{textcomp}%
\usepackage{manyfoot}%
\usepackage{booktabs}%
\usepackage{algorithm}%
\usepackage{algorithmicx}%
\usepackage{algpseudocode}%
\usepackage{listings}%
%%%%

%%%%%=============================================================================%%%%
%%%%  Remarks: This template is provided to aid authors with the preparation
%%%%  of original research articles intended for submission to journals published 
%%%%  by Springer Nature. The guidance has been prepared in partnership with 
%%%%  production teams to conform to Springer Nature technical requirements. 
%%%%  Editorial and presentation requirements differ among journal portfolios and 
%%%%  research disciplines. You may find sections in this template are irrelevant 
%%%%  to your work and are empowered to omit any such section if allowed by the 
%%%%  journal you intend to submit to. The submission guidelines and policies 
%%%%  of the journal take precedence. A detailed User Manual is available in the 
%%%%  template package for technical guidance.
%%%%%=============================================================================%%%%

%% as per the requirement new theorem styles can be included as shown below
\theoremstyle{thmstyleone}%
\newtheorem{theorem}{Theorem}%  meant for continuous numbers
%%\newtheorem{theorem}{Theorem}[section]% meant for sectionwise numbers
%% optional argument [theorem] produces theorem numbering sequence instead of independent numbers for Proposition
\newtheorem{proposition}[theorem]{Proposition}% 
%%\newtheorem{proposition}{Proposition}% to get separate numbers for theorem and proposition etc.

\theoremstyle{thmstyletwo}%
\newtheorem{example}{Example}%
\newtheorem{remark}{Remark}%

\theoremstyle{thmstylethree}%
\newtheorem{definition}{Definition}%

\raggedbottom
%%\unnumbered% uncomment this for unnumbered level heads

\begin{document}

\title[Identifying consensus tumor cell states and characterizing their roles in lung adenocarcinoma progression and clinical outcomes]{Identifying consensus tumor cell states and characterizing their roles in lung adenocarcinoma progression and clinical outcomes}

%%=============================================================%%
%% GivenName	-> \fnm{Joergen W.}
%% Particle	-> \spfx{van der} -> surname prefix
%% FamilyName	-> \sur{Ploeg}
%% Suffix	-> \sfx{IV}
%% \author*[1,2]{\fnm{Joergen W.} \spfx{van der} \sur{Ploeg} 
%%  \sfx{IV}}\email{iauthor@gmail.com}
%%=============================================================%%

\author[1]{\fnm{Shaobo} \sur{Kang}}
\equalcont{These authors contributed equally to this work as first authors.}

\author[1]{\fnm{Yijing} \sur{Zhang}}
\equalcont{These authors contributed equally to this work as first authors.}

\author[1]{\fnm{Renjie} \sur{Dou}}
\equalcont{These authors contributed equally to this work as first authors.}

\author*[1]{\fnm{Yanyan} \sur{Ping}}\email{pingyanyan@hrbmu.edu.cn}

\affil*[1]{
\orgdiv{College of Bioinformatics Science and Technology},
\orgname{Harbin Medical University},
\orgaddress{
\city{Harbin},
\postcode{150081},
\country{China}
}
}

%%==================================%%
%% Sample for unstructured abstract %%
%%==================================%%

\abstract{High heterogeneity of malignant cells across lung adenocarcinoma (LUAD) patients highlights the need to define conserved tumor cell states. Here, we identified nine consensus meta-programs (MPs) from single-cell transcriptomes of 30,142 malignant cells across 45 LUAD samples. Based on dominant MP activities, malignant cells were classified into nine consensus tumor cell states with distinct functional, regulatory, and microenvironmental features. The Cell-cycle, EMT, and Glycolysis states showed stronger outgoing communication with tumor-associated immune and stromal cells, and representative state-associated ligand–receptor interactions were linked to poor prognosis. Multi-resolution transcriptomic analyses further revealed progression-associated MP-state dynamics, with TNF\(\alpha\)--NF\(\kappa\)B and MHC class II states enriched in early-stage LUAD and Cell-cycle, EMT, and Glycolysis states more abundant in advanced-stage LUAD. Trajectory analysis suggested divergent transitions from TNF\(\alpha\)--NF\(\kappa\)B, MHC class II, Translation, and ATP synthesis states toward proliferative and invasive/metabolic states, accompanied by dynamic changes in state-specific transcription factors (TFs). Leveraging these TFs, we constructed an eight-TF prognostic signature consisting of KLF4, FOSL2, ATF1, POU6F1, FOXM1, TFDP1, XBP1, and CEBPG, which robustly predicted patient survival across multiple independent cohorts. Together, our study defines conserved tumor cell states in LUAD and links them to microenvironmental crosstalk, disease progression, and clinical outcome, providing a framework for biomarker discovery and translational investigation.}

%%================================%%
%% Sample for structured abstract %%
%%================================%%

% \abstract{\textbf{Purpose:} The abstract serves both as a general introduction to the topic and as a brief, non-technical summary of the main results and their implications. The abstract must not include subheadings (unless expressly permitted in the journal's Instructions to Authors), equations or citations. As a guide the abstract should not exceed 200 words. Most journals do not set a hard limit however authors are advised to check the author instructions for the journal they are submitting to.
% 
% \textbf{Methods:} The abstract serves both as a general introduction to the topic and as a brief, non-technical summary of the main results and their implications. The abstract must not include subheadings (unless expressly permitted in the journal's Instructions to Authors), equations or citations. As a guide the abstract should not exceed 200 words. Most journals do not set a hard limit however authors are advised to check the author instructions for the journal they are submitting to.
% 
% \textbf{Results:} The abstract serves both as a general introduction to the topic and as a brief, non-technical summary of the main results and their implications. The abstract must not include subheadings (unless expressly permitted in the journal's Instructions to Authors), equations or citations. As a guide the abstract should not exceed 200 words. Most journals do not set a hard limit however authors are advised to check the author instructions for the journal they are submitting to.
% 
% \textbf{Conclusion:} The abstract serves both as a general introduction to the topic and as a brief, non-technical summary of the main results and their implications. The abstract must not include subheadings (unless expressly permitted in the journal's Instructions to Authors), equations or citations. As a guide the abstract should not exceed 200 words. Most journals do not set a hard limit however authors are advised to check the author instructions for the journal they are submitting to.}

\keywords{Single-cell transcriptomics, Consensus tumor cell states, Transcription factors}

%%\pacs[JEL Classification]{D8, H51}

%%\pacs[MSC Classification]{35A01, 65L10, 65L12, 65L20, 65L70}

\maketitle
\section{Introduction}\label{sec1}
Lung adenocarcinoma (LUAD) is a common malignant tumor originating from glandular epithelial cells in lung tissue \cite {bib1}. Malignant epithelial cells in LUAD patients exhibit high heterogeneity, which is one of the primary factors contributing to the failure of cancer treatments \cite{bib2}. It is essential to define recurrent tumor cell states for understanding the common biological principles underlying LUAD progression and for identifying clinically relevant biomarkers.

Recent single-cell transcriptomic studies have begun to delineate malignant-cell heterogeneity in LUAD. Clustering-based analyses have identified distinct malignant epithelial subpopulations, including lineage-related, proliferative, and stem-like tumor cell states \cite{bib3,bib4}. Studies of early-stage LUAD have further revealed tumor clusters enriched in translation initiation and antigen presentation, as well as dynamic transitions between immune-responsive states and cell death-related states during the progression from pre-invasive to invasive lesions \cite{bib5,bib6}. Additional trajectory and metastasis‑oriented analyses indicated that malignant cells can acquire transcriptional states linked to either disease progression or aggressive phenotypes \cite{bib7,bib8}. Moreover, KRT8+ alveolar intermediate cell states with increased plasticity have been reported in early LUAD and related premalignant contexts \cite{bib9,bib10}. However, these studies defined tumor cell states through clustering-based annotation after integrating multiple samples. This strategy may be influenced by inter-patient variation and may not directly identify recurrent tumor cell states shared across patients.

Non-negative matrix factorization (NMF)-based approaches provide an effective framework for addressing this question. By decomposing malignant-cell transcriptomes within individual tumors, NMF can identify transcriptional programs that reflect intratumoral heterogeneity. When robust programs are subsequently compared and integrated across patients, recurrent meta-programs (MPs) can be identified as cross-sample. Previous studies have identified recurrent malignant-cell programs such as partial EMT programs in head-and-neck cancer and cell-cycle, interferon-response, stress-response, EMT, and metabolic programs across multiple cancer types \cite{bib11,bib12}. Cancer-type-specific analyses have also revealed consensus programs in nasopharyngeal carcinoma and lung cancer contexts \cite{bib13,bib14,bib15}. More recently, Gavish et al. identified 41 consensus meta-programs from 1,163 tumor samples spanning 24 cancer types, highlighting conserved transcriptional principles underlying malignant-cell heterogeneity \cite{bib16}. However, a systematic characterization of MP-defined conserved tumor cell states in LUAD, together with their functional, regulatory, microenvironmental, and clinical relevance, is still lacking.

Here, we bridge this gap by identifying nine consensus meta-programs across LUAD patients using NMF. We characterized the tumor cell states defined by these MP activities in terms of function, intercellular interaction, developmental trajectory, transcriptional regulation, and prognostic implications. This work systematically characterizes conserved tumor cell states underlying LUAD heterogeneity and provides a framework for understanding their links to microenvironmental crosstalk, tumor progression, and clinical outcome.

\section{Results}\label{sec2}

\subsection{Identification of malignant epithelial cells in LUAD patients}\label{subsec1}

We collected six single-cell RNA datasets containing 432,276 cells from 88 samples (20 normal, 68 tumor). The batch effect between different samples was corrected using Harmony package (Supplementary Fig. 1A). After initial quality control, we annotated eight major cell types (including epithelial cells, T cells, NK cells, B cells, myeloid cells, mast cells, fibroblast cells, and endothelial cells) through unsupervised clustering (Supplementary Fig. 1B), and canonical markers (Fig.~\ref{Figure1.png}A). Those clusters exhibiting marker genes from multiple distinct cell lineages were further excluded (Supplementary Fig. 1C). The proportions of each cell type varied across the LUAD patients (Fig.~\ref{Figure1.png}B). Additionally, we annotated the subtypes of the main cell types within the tumor microenvironment to further investigate its composition (Supplementary Fig. 2).
To identify malignant epithelial cells in LUAD patients, we re-clustered 30,142 epithelial cells into 26 cell clusters after removing potential doublet clusters (Supplementary Fig. 1D,E). We define the malignant cell clusters based on the assumption that the vast majority of cells in the cluster originate from tumor tissue (Fig.~\ref{Figure1.png}C). To further verify the reliability of malignant and non-malignant clusters, we used InferCNV to infer large-scale chromosomal copy-number variations (CNVs) of malignant and non-malignant cell clusters from tumor tissues by taking the transcriptome of the normal cells from non-tumor tissue as reference. We found that the identified malignant clusters showed obvious copy number alterations across the genome. However, the copy number profiles of the identified non-malignant clusters from tumor tissues were similar to that of reference cells (Fig.~\ref{Figure1.png}D). The identified malignant clusters showed higher CNV scores than normal clusters and high expression of LUAD markers MDK and TIMP1 (Fig.~\ref{Figure1.png}E-G) \cite{bib4}. These results suggested the reliability of the identified malignant cell clusters.

\subsection{Identification of nine consensus meta-programs}\label{subsec2}

Given the observed transcriptional heterogeneity of malignant cells across LUAD patients (Supplementary Fig. 3A), we identified nine consensus meta-programs (MPs) across patients, and each MP was represented by the top 50 most frequent genes within its cluster ( Fig.~\ref{Figure2.png}A, Supplementary Table. 3).

Functional enrichment analysis demonstrated that these MPs exhibited distinct functions (Fig.~\ref{Figure2.png}B). Both MP\_1 and MP\_2 were strongly associated with cell cycle and proliferation pathways: MP\_1 was significantly enriched in functions of G2M checkpoint, mitotic cell cycle, and cell division, whereas MP\_2 was predominantly characterized by DNA replication and DNA metabolic processes. Both MP\_8 and MP\_9 exhibited signatures associated with hypoxia and glycolysis. MP\_8 was specifically enriched in pathways related to epithelial-mesenchymal transition (EMT), IL2/STAT5 signaling, and KRAS activation, indicating a pro-metastatic program. In contrast, MP\_9 was dominated by mTORC1 signaling, indicating a pro-proliferation program.

Most genes in MP\_3 were involved in TNF-$\alpha$ signaling via NF-$\kappa$B, a central pathway that mediates inflammation, cell survival, and resistance to apoptosis in the tumor microenvironment. Meanwhile, MP\_3 was also associated with stress response pathways, including UV response and response to oxygen-containing compounds. These functional enrichments indicate that MP\_3 represented a pro-inflammatory and stress-responsive program. MP\_4 represented an immunogenic program associated with anti-tumor immune recognition, as it was specifically enriched in MHC class II-mediated antigen presentation pathways, including exogenous antigen processing and MHC-peptide complex assembly.
MP\_5 represented a translation-related program, specifically enriched in translation-associated functions such as translation initiation. MP\_6 was an energy metabolism-related program, specifically enriched in oxidative phosphorylation and aerobic respiration pathways, supporting efficient ATP synthesis and energy production in tumor cells. MP\_7 was involved in cell adhesion and apoptotic process. The overlap with the previously reported 41 MPs validated that these nine MPs were conserved transcriptional programs (Supplementary Fig. 3B).

\subsection{Identification and transcriptional characterization of nine consensus tumor cell states dominated by the meta-programs}\label{subsec3}

We calculated the activity scores of the nine consensus meta-programs (MPs) in all malignant epithelial cells and retained cells with at least one activated MP. Notably, the activation patterns of these nine MPs exhibited strong mutual exclusivity across individual tumor cells, with each cell predominantly dominated by only one MP signature. Accordingly, these malignant cells were stratified into nine distinct tumor cell states based on their dominant MP profiles, and each state was named according to its corresponding MP (Fig.~\ref{Figure2.png}C). The state-specific high expression of MP signatures (Fig.~\ref{Figure2.png}D, Supplementary Fig. 3D) and high activities of corresponding hallmarks (Fig.~\ref{Figure2.png}E) supported the dominant role of each MP in defining its cell state. Meanwhile, these MP states were commonly observed across different LUAD patients, with multiple distinct states coexisting within the same tumor sample (Supplementary Fig. 3C), suggesting that these nine MP-defined tumor cell states are consistently present across individuals.

We further constructed state-specific functional profiles for the nine MP states using their highly and specifically expressed genes. Each state exhibited a distinct functional signature, with high activity in pathways corresponding to its dominant MP program (Fig.~\ref{Figure3.png}A, Supplementary Fig. 4A). For example, the Cell-cycle states (MP\_1 and MP\_2) showed divergent functional activities despite both being cell cycle-active: MP\_1 was strongly enriched in mitotic progression, whereas MP\_2 was more associated with DNA replication and DNA damage repair. Similarly, the EMT state (MP\_8) and Glycolysis state (MP\_9) also displayed distinct functional biases: MP\_8 was linked to cell adhesion and migration, whereas MP\_9 showed high activity in glycolytic and hypoxic metabolism pathways. These divergent functional profiles highlight the functional heterogeneity across the nine tumor cell states and further support the robustness and biological specificity of our state definitions.

Furthermore, we identified 56 state-specific transcription factors (TFs) across MP states using pySCENIC analysis (Fig.~\ref{Figure3.png}B). The unique transcriptional regulatory programs of each tumor cell state were demonstrated by the specific regulon activity of these TFs and their significant positive correlations with MP activity scores (Fig.~\ref{Figure3.png}B-C, Supplementary Fig. 4B-C). Meanwhile, we found that these state-specific TFs directly and specifically targeted MP signature genes (Fig.~\ref{Figure3.png}D). For example, the Cell-cycle states (MP\_1 and MP\_2) exhibited nearly distinct TF sets. VPS72 and FOXM1 showed specific activation and positive correlations with the MP\_1 state, whereas TFDP1 and E2F family members were specifically associated with the MP\_2 state.
Notably, VPS72 has rarely been studied in LUAD. We found that VPS72 directly and specifically targeted UBE2C, a gene reported in previous studies to mark a proliferative tumor subpopulation in LUAD. VPS72 showed specific high expression in malignant cells (Supplementary Fig. 5B) and was spatially confined to tumor regions (Supplementary Fig. 5C), a pattern further supported by bulk transcriptomic cohorts (Fig. S5C). Consistently, VPS72 expression was positively correlated with UBE2C across seven independent LUAD cohorts (Supplementary Fig. 5A). Furthermore, elevated VPS72 expression was significantly associated with poor prognosis (Supplementary Fig. 5D). These results suggest a potential regulatory role of VPS72 in the MP\_1 state.

\subsection{Identification of shared and state-specific intercellular communication patterns among nine consensus tumor cell states}\label{subsec4}

To explore the heterogeneity of crosstalk between the nine MP-defined tumor cell states and the tumor microenvironment, we performed CellChat analysis between MP states and tumor-associated immune and stromal subtypes that showed significant abundance differences between tumor and normal tissues (Supplementary Fig. 6). Compared with other tumor cell states, the Cell-cycle states (MP\_1 and MP\_2), Cell adhesion state (MP\_7), EMT state (MP\_8), and Glycolysis state (MP\_9) exhibited stronger outgoing signaling activities, whereas their incoming signaling patterns were broadly similar (Fig.~\ref{Figure4.png}A). These outgoing interactions were mainly directed toward SPP1+ macrophages, CXCL10+ macrophages, COL10A1+ fibroblasts, IGF1+ fibroblasts, and CCL21+ endothelial cells, suggesting that aggressive MP states may shape a pro-tumor microenvironment.

A total of 106 significant ligand–receptor pairs were identified and showed heterogeneous interaction strengths among MP states (Supplementary Fig. 7). Among them, MDK--NCL was a broadly activated ligand–receptor pair between malignant cell states and tumor-associated cell subsets, with markedly stronger interactions in MP\_1, MP\_2, MP\_7, MP\_8 and MP\_9 states. MDK expression was significantly elevated in these five MP states compared to other malignant subpopulations and NCL widely expressed in tumor-associated cell populations (Fig.~\ref{Figure4.png}B). Previous studies have reported that elevated MDK expression is associated with malignant epithelial phenotypes and that MDK--NCL signaling can promote endothelial cell migration and cancer-associated fibroblast proliferation \cite {bib17,bib18,bib19}. Spatial transcriptomic analysis further showed that MDK+NCL+ spots were predominantly enriched within cancer regions in the AIS, MIA and IAC samples (Fig.~\ref{Figure4.png}C). Moreover, patients with high expression of both MDK and NCL showed the poorest survival outcomes in four independent LUAD cohorts (Fig.~\ref{Figure4.png}D). Elevated MDK in these MP states amplified MDK–NCL crosstalk with tumor-associated cells and might remodel the local tumor microenvironment.
In contrast, SCGB3A2--MARCO showed stronger interaction strength between the MP\_3, MP\_4, MP\_5 and MP\_6 states and macrophage subpopulations. SCGB3A2 was preferentially expressed in the MP\_3, MP\_4, MP\_5 and MP\_6 states (Fig.~\ref{Figure4.png}E). Spatial analysis showed that SCGB3A2+MARCO+ spots were mainly enriched within tumor regions of early-stage lesions (AIS and MIA) (Fig.~\ref{Figure4.png}F). Patients with high expression of both SCGB3A2 and MARCO exhibited better survival (Fig.~\ref{Figure4.png}G).

We further explored state-specific ligand--receptor interactions (Supplementary Fig. 8). The interactions of ANGPTL4--CDH11 between the MP\_8 state and the endothelial cell subpopulations, and the interaction of ANGPTL4--CDH11 between the MP\_8 state and the fibroblast cell subpopulations were observed specially in the MP\_8 state, with ANGPTL4 selectively expressed in MP\_8 and MP\_9 states and CDH5/CDH11 preferentially expressed in endothelial and fibroblast cell subpopulations, respectively. In addition, PLAU--PLAUR interaction was also specific to MP\_8 state, with PLAUR mainly expressed in macrophage and fibroblast cell subpopulations. Spatial co-localization and survival analyses further supported the association of these state-specific interactions with tumor regions and poor prognosis.
These results indicated that MP\_8 state specifically secrete ANGPTL4 and PLAU to recruit endothelial, fibroblast and macrophage cell subpopulations for remodeling a migration‑favorable microenvironment.

\subsection{Multi-resolution characterization of progression-associated dynamics of nine consensus tumor cell states}\label{subsec5}

To investigate whether the nine consensus tumor cell states were associated with LUAD progression, we examined MP activity patterns and MP-state proportions across single-cell samples. The MP states showed clear stage-associated differences among LUAD samples (Fig.~\ref{Figure5.png}A). During progression from early to advanced LUAD, the proportions of the TNF-\(\alpha\)--NF-\(\kappa\)B state (MP\_3) and MHC class II state (MP\_4) were significantly decreased, whereas the proportions of the Cell-cycle states (MP\_1 and MP\_2), EMT state (MP\_8), and Glycolysis state (MP\_9) were significantly increased (Fig.~\ref{Figure5.png}B, Supplementary Fig. 9A).

To validate these progression-associated changes in bulk transcriptomic datasets, we evaluated MP-state signatures in TCGA-LUAD samples. Consistent with the single-cell results, ssGSEA analysis showed that the scores of MP\_3 and MP\_4 states were reduced during LUAD progression, whereas the scores of MP\_1, MP\_2, MP\_8, and MP\_9 states progressively increased from normal tissues to early-stage and advanced-stage LUAD. Deconvolution analyses using BayesPrism and CIBERSORT further confirmed these stage-associated trends (Fig.~\ref{Figure5.png}C and Supplementary Fig. 9B, C).
We next examined the spatial distribution of MP states using six LUAD spatial transcriptomics samples from GSE189487. The proportions of MP states within cancer spots were estimated using Robust Cell Type Decomposition (RCTD). Across these spatial samples, the MP\_4 state was mainly enriched in AIS lesions, whereas the MP\_1, MP\_2, and MP\_8 states were preferentially enriched in IAC lesions (Fig.~\ref{Figure5.png}D, E, Supplementary Fig. 9D, E). Collectively, these multi-resolution analyses revealed a robust association between MP-state dynamics and LUAD progression, providing a rationale for further investigation of potential transition relationships among the nine MP-defined tumor cell states.

\subsection{Divergent transcriptional trajectories and TF-regulatory dynamics among nine consensus tumor cell states}\label{subsec6}

Given the progression-associated changes of MP states observed in multi-resolution transcriptomic analyses, we next investigated the potential transition relationships among these MP-defined tumor cell states. Monocle3 trajectory analysis partitioned MP-state tumor cells into seven transcriptionally distinct clusters (Fig.~\ref{Figure6.png}A). Cluster 2, mainly composed of the TNF-$\alpha$--NF-$\kappa$B state (MP\_3), MHC class II state (MP\_4), Translation state (MP\_5), and ATP synthesis state (MP\_6), showed the lowest CNV and CytoTRACE-derived stemness scores and was therefore defined as the trajectory root (Fig.~\ref{Figure6.png}B, Supplementary Fig. 10A).

The reconstructed trajectory exhibited a clear branched structure. Lineage 1 extended from Cluster 2 toward Cluster 1 through an intermediate Cluster 4, accompanied by decreased proportions of the MHC class II state (MP\_4), Translation state (MP\_5), and ATP synthesis state (MP\_6) and increased proportions of the Cell adhesion state (MP\_7), EMT state (MP\_8), and Glycolysis state (MP\_9). This branch suggested a transition toward invasive and metabolically active tumor cell states. In contrast, Lineage 2 extended from Cluster 2 toward Cluster 3 through Cluster 6, where the Cell-cycle states (MP\_1 and MP\_2) gradually became dominant, indicating a transition toward proliferative tumor cell states. Lineage 3, ending in Cluster 5, showed an MP-state composition similar to Cluster 4, suggesting that it may represent an additional transitional branch rather than an independent terminal state (Fig.~\ref{Figure6.png}C). Slingshot analysis further supported the branched structure inferred by Monocle3 (Fig.~\ref{Figure6.png}D, Supplementary Fig. 10B). Together, these results suggest that LUAD tumor cells may diverge from relatively low-CNV and low-stemness states toward either invasive/metabolic or proliferative trajectories.

We further examined TF-regulatory dynamics along the two major trajectory branches. TFs associated with the TNF-$\alpha$--NF-$\kappa$B state (MP\_3) and MHC class II state (MP\_4), including CEBPA, XBP1, ETV1, and FOXA2, showed gradually decreased activities along both lineages. In contrast, lineage-specific TF activation patterns were observed in the two terminal directions. Along Lineage 1, EMT-associated TFs were progressively activated, with FOSL2 showing a marked increase along pseudotime. Along Lineage 2, cell-cycle-associated TFs were progressively activated, with FOXM1 and VPS72 showing increased activities toward the proliferative branch (Fig.~\ref{Figure6.png}E, Supplementary Fig. 10C). GeneSwitches analysis further supported the activation of FOSL2 along the invasive trajectory and VPS72 along the proliferative trajectory (Fig.~\ref{Figure6.png}F). Collectively, these findings suggest that MP-state transitions are associated with distinct TF-regulatory programs, which may contribute to invasive and proliferative tumor cell trajectories during LUAD progression.

\subsection{Clinical relevance of MP-state-specific TFs in LUAD prognosis and therapeutic response}\label{subsec7}

Given the progression-associated dynamics and TF-regulatory trajectories of MP-defined tumor cell states, we next evaluated their clinical relevance in LUAD. By integrating TCGA-LUAD survival data with malignant single-cell transcriptomes using Scissor, we found that poor-prognosis-associated MP states were mainly enriched in the Cell-cycle states (MP\_1 and MP\_2), Cell adhesion state (MP\_7), EMT state (MP\_8), and Glycolysis state (MP\_9), whereas favorable-prognosis-associated MP states were dominated by the TNF-$\alpha$--NF-$\kappa$B state (MP\_3) and MHC class II state (MP\_4) (Supplementary Fig. 11A). This pattern was consistently supported across seven independent validation cohorts, indicating that distinct MP states are associated with divergent clinical outcomes (Fig.~\ref{Figure7.png}A).

We then examined whether MP-state-specific TFs could be translated into a prognostic model. Univariate Cox regression identified 10 candidate TFs associated with overall survival (Supplementary Fig. 12A), which were further evaluated using 117 integrative machine-learning models (Supplementary Fig. 11B). Among these models, the StepCox (forward) + Enet ($\alpha = 0.7$) model showed the best overall performance and generated an eight-TF prognostic signature consisting of ATF1, CEBPG, FOSL2, FOXM1, KLF4, POU6F1, TFDP1, and XBP1 (Fig.~\ref{Figure7.png}B). This model showed stable predictive performance, with a mean 1-year AUC of 0.758 across seven validation cohorts (Supplementary Fig. 12C).
Based on this eight-TF signature, we calculated a risk score for each LUAD patient. The risk score significantly stratified patients into high- and low-risk groups in the TCGA-LUAD cohort, and this prognostic stratification was consistently validated across seven independent cohorts (Fig.~\ref{Figure7.png}C). Multivariate Cox regression further showed that the risk score remained an independent prognostic factor after adjustment for age and clinical stage (HR = 3.53; Supplementary Fig. 12B). Significant survival differences between risk groups were also observed within different clinical stages, supporting the stability of this TF-based model across LUAD patients with distinct disease stages (Fig.~\ref{Figure7.png}D). Consistently, the high-risk group contained a higher proportion of advanced-stage patients (Supplementary Fig. 12D).

We further assessed the biological and therapeutic relevance of the TF-based risk score. ESTIMATE and immune infiltration analyses showed that the high-risk group was associated with higher tumor purity and increased macrophage infiltration, whereas the low-risk group showed higher immune and stromal scores and greater enrichment of na\"ive B cells and CD8$^+$ T cells (Supplementary Fig. 12E, F). Drug sensitivity analysis identified 96 compounds significantly correlated with the risk score, with low-risk patients showing greater predicted sensitivity to nutlin-3, nilotinib, bexarotene, and MK-2206, and high-risk patients showing greater predicted sensitivity to MK-1775 and paclitaxel (Fig.~\ref{Figure7.png}E, F). Collectively, these results indicate that MP-state-specific TFs are closely associated with LUAD prognosis, immune microenvironmental features, and predicted therapeutic sensitivity, highlighting their potential clinical relevance.


%%=============================================%%
%% For presentation purpose, we have included  %%
%% \bigskip command. Please ignore this.       %%
%%=============================================%%

\begin{figure}[h]
\centering
\includegraphics[height=0.75\textheight, keepaspectratio]{Figure1.png}
\caption{
A single-cell atlas of normal lung tissues and LUAD tumor tissues.
\textbf{A}. Uniform Manifold Approximation and Projection (UMAP) plots showing the major cell types identified from 432,276 cells across 88 LUAD specimens (left). Bubble heatmap showing the expression of canonical marker genes for the major cell types (right). Dot size represents the proportion of cells expressing each marker gene, and color intensity indicates the mean expression level.
\textbf{B}. Relative proportions of major cell types in each sample.
\textbf{C}. Relative proportions of tissue origins (top) and sample origins (bottom) for each epithelial cell cluster.
\textbf{D}. Heatmap showing large-scale copy number variations (CNVs) in epithelial cells. Normal cells from non-tumor tissues are shown in the upper heatmap, whereas normal and tumor epithelial cells from tumor tissues are shown in the lower heatmap. Genes are ordered from left to right according to their genomic positions across chromosomes. Red indicates amplifications, and blue indicates deletions.
\textbf{E}. Marker gene expression patterns of normal epithelial cell subpopulations and tumor cells.
\textbf{F}. Boxplot showing CNV scores of epithelial cells from normal and tumor tissues.
\textbf{G}. UMAP plot of 30,142 epithelial cells, color-coded by cell subset and CNV score.
}
\label{Figure1.png}
\end{figure}

\begin{figure}[h]
\centering
\includegraphics[width=0.9\textwidth]{Figure2.png}
\caption{
Nine tumor cell state with consensus expression programs among patients.
\textbf{A}. Heatmap showing Jaccard similarity indices among 185 robust NMF programs derived from 45 LUAD samples. Programs were hierarchically clustered and grouped into nine meta‑programs (MPs).
\textbf{B}. Functional annotation of the nine MPs based on Gene Ontology (GO) enrichment analysis.
\textbf{C}. Heatmap shows activation scores of MPs in MP states.
\textbf{D}. Heatmap depicts the relative signature expressions of MPs among MP states.
\textbf{E}.  Violin plots showing the specificity of pathway activities among different MP states.
}
\label{Figure2.png}
\end{figure}

\begin{figure}[h]
\centering
\includegraphics[width=0.9\textwidth]{Figure3.png}
\caption{
Transcription factor activities and regulatory networks shaping MP states.
\textbf{A}. GO enrichment was performed for upregulated genes of each MP state and visualized as a network. Each node represents a GO term, sized by gene count, and edges indicate functional similarity (Jaccard index). Nodes are colored by MP states were grouped into major biological modules.
\textbf{B}. Bubble diagram showing specificity TF-regulon for each MP state. Dot size indicates regulon specificity score (RSS), and color represents regulon activity (Z-score).
\textbf{C}. Pearson correlation between TF regulon activity and MP activity scores.
\textbf{D}. TF-meta-program gene regulatory network. V-shaped nodes represent TFs, while circular nodes represent genes. Edges indicate regulatory relationships between TFs and genes, with line thickness proportional to interaction strength. The nodes and edges are colored according to their MP state.
}
\label{Figure3.png}
\end{figure}

\begin{figure}[h]
\centering
\includegraphics[height=0.75\textheight, keepaspectratio]{Figure4.png}
\caption{
Shared and state-specific intercellular communication patterns of MP states within the tumor microenvironment.
\textbf{A}. Dot plot show the overview of signal sending strengths from MP states to tumor-associated cells (left) and signal reception strengths from tumor-associated cells to MP states (right).
\textbf{B,E}. Boxplots displaying the strength of representative ligand–receptor (L–R) pairs across MP states. Violin plots show expression distributions of corresponding L-R genes across cell types.
\textbf{C,F}. Spatial distribution of selected L–R interaction groups in six spatial transcriptomics samples, highlighting their localization in tumor regions. AIS, sample TD5 and TD8; MIA, sample TD3 and TD6; IAC, sample TD1 and TD2.
\textbf{D,G}. Kaplan–Meier survival analyses of LUAD patients stratified by MDK–NCL (D) and SCGB3A2-MARCO (G) interaction-related expression patterns across independent cohorts.
}
\label{Figure4.png}
\end{figure}

\begin{figure}[h]
\centering
\includegraphics[width=0.9\textwidth]{Figure5.png}
\caption{
Nine MP states associated with the progression of LUAD.
\textbf{A}. Heatmap of mean MP activity scores in tumor cells across single-cell LUAD samples, Color represents MP score (blue: lower, red: higher).
\textbf{B}. Boxplots showing cell ratio differences of MP states between Early (blue) and Advanced (red) LUAD stages in single-cell data.
\textbf{C}. The scores of MP state signature were calculated by ssGSEA (top); Proportion of the nine MP states estimated by BayesPrism for TCGA LUAD samples at different stages (bottom). P.adjust.method = wilcox.test.
\textbf{D}.  Representative spatial transcriptomic maps showing the distribution of MP states in LUAD samples with different pathological progression stages. Sample ST8 from AIS (Adenocarcinoma in Situ), ST3 from MIA (Minimally Invasive Adenocarcinoma), and ST1 from IAC (Invasive Adenocarcinoma).
\textbf{E}. Boxplots quantifying cell ratios of MP states in cancer spots.
}
\label{Figure5.png}
\end{figure}

\begin{figure}[h]
\centering
\includegraphics[width=0.9\textwidth]{Figure6.png}
\caption{
Transcriptional trajectory analysis of between MP states.
\textbf{A}. Unsupervised transcriptional trajectory of the MP states, colored by Monocle3-derived clusters (left) and stage (right). MP states were partitioned into seven transcriptionally distinct clusters.
\textbf{B}. Violin plots showing CNV scores (top) and stemness scores (bottom) across the seven clusters.
\textbf{C}. Monocle3 pseudotime trajectory projected onto the UMAP embedding. Pie charts indicate the composition of MP states within each cluster. 
\textbf{D}. Slingshot-based trajectory inference validating the branching structure. Three inferred lineages are shown on the UMAP embedding.
\textbf{E}. Dynamic changes in TF regulon activities along pseudotime in different lineages. Heatmap showing TF activity patterns across pseudotime, with representative TFs associated with lineage-specific state transitions highlighted on the right.
\textbf{F}. GeneSwitches analysis of gene expression dynamics along pseudotime. Representative TFs and MP genes are shown for Lineage 1 and Lineage 2.
}
\label{Figure6.png}
\end{figure}

\begin{figure}[h]
\centering
\includegraphics[width=0.9\textwidth]{Figure7.png}
\caption{
The prognostic effect of nine MP states.
\textbf{A}. The ratio of different tumor states in Scissor+ cells and Scissor- cells in eight independent datasets.
\textbf{B}. Prognostic values of the MP state- specific TFs in eight bulk cohorts by univariate cox regression analysis. TF colors represent their corresponding MP states.
\textbf{C}. Survival analysis in TCGA and seven independent cohorts. Calculation of the risk score from the product of the expression of the TFs and the risk coefficients.
\textbf{D}. Kaplan–Meier survival analysis of high- and low-risk groups within stage I, II, III, and IV LUAD patients.
\textbf{E}. Volcano plot showing correlations between RS and predicted drug sensitivity (IC50). Compounds significantly positively or negatively correlated with RS are highlighted (Spearman correlation \(>0.3\) or \(<-0.3\), adjusted \(P<0.05\)).
\textbf{F}. Representative correlations between RS and predicted drug sensitivity (IC50) for selected compounds.
}
\label{Figure7.png}
\end{figure}

\section{Methods}\label{sec11}

\subsection{Materials}

\subsubsection{Single-cell RNA-seq expression profiles}

We downloaded five LUAD single-cell RNA-seq datasets from the Curated Cancer Cell Atlas (3CA; \url{https://www.weizmann.ac.il/sites/3CA}), including datasets from Bischoff et al., Kim et al., Laughney et al., Qian et al., and Xing et al. \cite{bib16}. In addition, one LUAD single-cell RNA-seq dataset was obtained from the Gene Expression Omnibus (GEO; GSE189357) \cite{bib3}. LUAD patients were divided into early-stage LUAD, including minimally invasive adenocarcinoma (MIA) and stage I tumors, and advanced-stage LUAD, including stage II--III tumors. These six datasets were used for subsequent analyses and included 20 adjacent non-tumor lung samples, 44 early-stage LUAD samples, and 24 advanced-stage LUAD samples (Supplementary Table. 1).

\subsubsection{Bulk transcriptomic and clinical data of LUAD}

We analyzed one RNA-seq dataset from The Cancer Genome Atlas (TCGA-LUAD) and seven microarray datasets from GEO, including GSE13213 \cite{bib22}, GSE31210 \cite{bib23}, GSE41271 \cite{bib24}, GSE26939 \cite{bib25}, GSE30219 \cite{bib26}, GSE72094 \cite{bib27}, and GSE11969 \cite{bib28}. All datasets included matched clinical information and were used for bulk-level validation and survival analyses (Supplementary Table. 2).

\subsubsection{Spatial transcriptomics data of LUAD}

LUAD spatial transcriptomics data were obtained from GEO dataset GSE189487 \cite{bib3}. This dataset included two adenocarcinoma in situ (AIS) samples (TD5 and TD8), two minimally invasive adenocarcinoma (MIA) samples (TD3 and TD6), and two invasive adenocarcinoma (IAC) samples (TD1 and TD2). Pathologically annotated hematoxylin and eosin (H\&E)-stained cancer, normal epithelial, stromal, and lymphatic subregions provided by Zhu et al. were reproduced through dimensionality reduction, clustering, spatial variable feature detection, and interactive visualization using Seurat.

\subsection{Ethics statement}

Not applicable.

\subsection{Methods}

\subsubsection{scRNA-seq data processing}

Single-cell RNA-seq data were processed using the Seurat package (v5) \cite{bib30}. Ambient RNA contamination was detected using the \texttt{decontX} R package, and only cells with a contamination score below 0.2 were retained \cite{bib31}. Quality control filtering was then applied to select high-quality cells according to the following criteria: (i) the proportion of mitochondrial genes was less than 10\%; (ii) the proportion of hemoglobin genes was less than 1\%; and (iii) the number of detected genes per cell (\texttt{nFeature\_RNA}) ranged from 500 to 7000. Genes detected in fewer than ten cells were excluded, as were mitochondrial, hemoglobin, and ribosomal genes. Potential doublets were identified and removed using the \texttt{scDblFinder} R package \cite{bib32}. In addition, clusters that co-expressed canonical markers of multiple distinct cell lineages were excluded to ensure population purity. Finally, a total of 296,429 single-cell transcriptomes were retained after quality control.

\subsubsection{Cell annotation}

Gene expression data were normalized and scaled using the \texttt{NormalizeData} and \texttt{ScaleData} functions, respectively. The top 2,000 highly variable genes were identified using \texttt{FindVariableFeatures}. Batch effects across patients were corrected using the Harmony package. Principal component analysis (PCA) was performed on the scaled data, and the first 30 principal components were selected for subsequent clustering and visualization. Cell clustering was performed using the \texttt{FindClusters} function at a resolution of 0.5, and uniform manifold approximation and projection (UMAP) was applied for two-dimensional visualization using \texttt{RunUMAP}. Clusters were annotated based on canonical marker genes \cite{bib29}. Clusters exhibiting marker genes from multiple distinct cell lineages were excluded from downstream analyses.

\subsubsection{Identification of malignant cells}

Epithelial cells were re-clustered, and malignant cell clusters were identified based on the criterion that the majority of cells within a cluster originated from tumor tissues. To further validate malignant identity, we examined the expression of LUAD marker genes, including \textit{MDK} and \textit{TIMP1}, and performed copy number variation (CNV) inference using the inferCNV package \cite{bib33}. Specifically, inferCNV was applied to infer large-scale chromosomal CNVs in malignant and non-malignant epithelial cells from tumor tissues, using normal epithelial cells from non-tumor tissues as references. The parameters were set as \texttt{cutoff = 0.1}, \texttt{denoise = TRUE}, and \texttt{scale\_data = TRUE}. The CNV score of each tumor cell was calculated by normalizing the inferred copy number values from inferCNV to a range of $-1$ to 1 and then calculating the sum of squared values.

\subsubsection{Identification of consensus meta-programs in malignant cells across patients}

Patients with more than 100 malignant cells were retained for subsequent analyses. To capture intratumoral transcriptional heterogeneity, non-negative matrix factorization (NMF) was performed separately for each sample. Multiple rank parameters, ranging from $K = 4$ to $K = 9$, were used, resulting in 39 programs for each tumor. Each program was represented by its top 50 most highly weighted genes.

To select robust and non-redundant programs, the following criteria were applied: (i) robustness within a tumor, where a program was required to share at least 70\% of its top genes with another program derived from the same tumor but with a different $K$ value; (ii) robustness across tumors, where a program was required to have at least 20\% gene overlap with any program from another tumor; and (iii) non-redundancy within a tumor, where programs were ranked by their maximal cross-tumor overlap and then selected sequentially, discarding any later program with more than 20\% overlap with an already selected program.

To consolidate robust NMF programs into meta-programs (MPs), all NMF programs were prioritized according to the number of programs from other tumors with which they shared at least 20\% gene overlap. The top-ranked program was used to seed a new cluster, and the initial MP was defined as its top 50 genes. The program with the highest gene overlap relative to the current MP was then added to the cluster, and the MP was updated by selecting the 50 genes that appeared most frequently among all programs in that cluster. This process was repeated until no additional program achieved at least 20\% overlap with the MP. The resulting MP was recorded, and the procedure was re-initiated on the remaining unclustered programs. Clusters containing fewer than five programs were discarded. Each MP comprised a gene set of 50 genes and was functionally annotated using Gene Ontology biological process gene sets (C5: GO BP) and Hallmark gene sets.

The implementation of NMF and consensus MP generation was adapted from the
\href{https://github.com/ZedekiahZhou/PanCan_scMetab}{PanCan\_scMetab workflow}
\cite{bib34}, which was modified from the
\href{https://github.com/tiroshlab/3ca/tree/main/ITH_hallmarks/Generating_MPs}
{\texttt{ITH\_hallmarks/Generating\_MPs} scripts}
developed by Gavish et al. \cite{bib16}.

\subsubsection{Assignment of tumor cell states based on MP activity}

To assign tumor cell states, the activity scores of each MP in malignant cells were calculated using the \texttt{sigScores} function from the scalop workflow (\url{https://github.com/jlaffy/scalop}). An MP was considered activated in a tumor cell if its activity score exceeded 0.8. For tumor cells with multiple activated MPs, the MP with the highest activity score was designated as the dominant program, and the corresponding cells were defined as being in that MP-defined tumor cell state. Cells without any activated MPs were classified as inactive tumor cells. To further validate the accuracy of MP-state assignment, pathway activity scores were calculated for each tumor cell using the \texttt{AddModuleScore} function.

\subsubsection{Differential expression and functional enrichment analyses}

MP-state-specific upregulated genes were identified using the \texttt{FindMarkers} function. Differentially expressed genes were required to be expressed in at least 25\% of cells in the corresponding cell state (\texttt{min.pct = 0.25}) and to have a log fold change greater than 1 (\texttt{logfc.threshold = 1}). Functional enrichment analysis of these differentially expressed genes was performed using the \texttt{compareCluster} function from the clusterProfiler R package. Pathway similarity was calculated using the Jaccard index with the \texttt{pairwise\_termsim} function (\texttt{method = "JC"}) and visualized using the \texttt{emapplot} function from the enrichplot R package with \texttt{clusterFunction = cluster::clara}.

\subsubsection{Identification of MP-state-specific transcription factors}

To investigate transcriptional regulation underlying MP-defined tumor cell states, pySCENIC \cite{bib35} was applied with default parameters to infer regulon activity scores using the count matrix as input. Based on the regulon activity scores, transcription factor (TF) specificity scores across different MP states were calculated using the \texttt{calcRSS} function. MP-state-specific TFs were selected using the \texttt{plotRSS} function with \texttt{thr = 0.01} and \texttt{zThreshold = 1.5}. To further assess the association between TFs and MP states, Pearson correlation coefficients were calculated between TF regulon activity scores and MP activity scores across tumor cells. Regulatory relationships between TFs and MP genes were visualized using Cytoscape \cite{bib36}, based on TF-target gene weights inferred by pySCENIC.

\subsubsection{Trajectory analysis and gene switching}

Dimensionality reduction and clustering were performed on all MP-state tumor cells using the standard Seurat workflow described above. The resulting Seurat object was converted into a Monocle3 object using the \texttt{new\_cell\_data\_set} function. Trajectory inference was performed using Monocle3 (v1.3.7) \cite{bib37} with default parameters. Specifically, the \texttt{cluster\_cells} and \texttt{learn\_graph} functions were used to re-cluster MP-state tumor cells and reconstruct the transcriptional trajectory, respectively. Tumor stemness scores were calculated using CytoTRACE \cite{bib38} to assist in determining the starting point of the trajectory. After defining the trajectory root, the \texttt{order\_cells} function was used to infer pseudotime for tumor cells along the trajectory. Based on the TF regulon activity matrix inferred by pySCENIC for MP-state cells in lineage 1 and lineage 2, diffusion maps were reconstructed using the ClusterGVis package.

Gene switching along pseudotime was inferred using the GeneSwitches package \cite{bib39}. Expression values were binarized with a cutoff of 0.2, and switching points were identified using logistic regression with the \texttt{find\_switch\_logistic\_fastglm} function. TFs and MP-associated genes with significant switching behavior were selected using the \texttt{filter\_switchgenes} function and were defined as candidate regulators of trajectory transitions.

\subsubsection{Cell--cell communication analysis}

Cell--cell communication analysis was performed using the CellChat package based on the normalized count matrix \cite{bib40}. Significant ligand--receptor (L--R) interactions between MP-defined malignant cell states and tumor-associated cells were inferred. The full \texttt{CellChatDB.human} database was used to infer ligand--receptor communications. Overexpressed ligands, receptors, and interactions were identified using \texttt{identifyOverExpressedGenes} and \texttt{identifyOverExpressedInteractions}. Communication probabilities for L--R pairs were computed using \texttt{computeCommunProb}, followed by quantification of interaction numbers and communication strength using \texttt{aggregateNet}.

\subsubsection{Estimation of MP-state proportions in bulk samples}

The reference matrix was first constructed using the \texttt{generateRef\_seurat} function from the IOBR R package \cite{bib41}, with the Seurat object normalized by \texttt{NormalizeData} as input. Based on this reference matrix, CIBERSORT was applied through the \texttt{deconvo\_tme} function with \texttt{method = "cibersort"} to estimate the relative proportions of distinct MP-defined tumor cell states in TCGA-LUAD samples.

BayesPrism was additionally applied as an orthogonal deconvolution approach \cite{bib42},. BayesPrism directly uses the single-cell count matrix restricted to protein-coding genes as input. The \texttt{run.prism} function was used to estimate the relative proportions of each MP state in TCGA-LUAD samples, and the resulting fractions for each bulk sample were obtained using the \texttt{get.fraction} function.

\subsubsection{Deconvolution analysis of spatial transcriptomics data}

To investigate the spatial distribution of distinct MP-defined tumor cell states, Robust Cell Type Decomposition (RCTD) \cite{bib43} was applied to deconvolute LUAD spatial transcriptomics data. Specifically, an RCTD object was constructed using the \texttt{create.RCTD} function from the spacexr R package, with the spatial spot count matrix, spatial coordinate information, and single-cell reference count matrix as input. The \texttt{run.RCTD} function was then used to estimate the composition of each MP-defined tumor cell state across spatial spots.

\subsubsection{Identification of prognosis-associated malignant cells using Scissor}

Scissor was used to identify phenotype-associated malignant cell subpopulations from single-cell data \cite{bib44}. LUAD overall survival data were integrated as the phenotype, together with the single-cell count matrix of malignant cells. Scissor was then applied to identify significantly positive prognosis-associated malignant cells (Scissor+) and significantly negative prognosis-associated malignant cells (Scissor-) using \(\alpha = 0.05\).

\subsubsection{Survival analysis}\label{subsubsec_survival}

The Mime1 R package was used to develop prognostic models from bulk transcriptomic data using 10 machine-learning algorithms \cite{bib45}. Prognosis-associated transcription factors (TFs) were first identified using the \texttt{ML.Dev.Prog.Sig} function with the parameter \texttt{unicox.filter.for.candi = TRUE}, yielding 10 candidate TFs. Subsequently, 117 integrative machine-learning models were constructed based on these candidate TFs, and the optimal model was selected according to the concordance index (C-index). The StepCox (forward) + Enet (\(\alpha = 0.7\)) model showed the best performance and was selected as the final model.

The risk score for each patient was calculated as follows:

\begin{equation}
\mathrm{Risk\ score} = \sum_{i=1}^{8}
\left( \beta_i \times \mathrm{Exp}_i \right),
\end{equation}

where \(\beta_i\) represents the coefficients of gene, and \(\mathrm{Exp}_i\) represents the expression level of the corresponding gene.

Patients were stratified into high- and low-risk groups according to the median risk score. Kaplan--Meier survival analysis was then performed to compare overall survival between the two groups. Model performance was evaluated using time-dependent receiver operating characteristic (ROC) analysis, and the area under the curve (AUC) values were calculated using the \texttt{cal\_AUC\_ml\_res} function. The prognostic relevance of VPS72 was further evaluated using the Kaplan--Meier Plotter database (\url{https://kmplot.com/analysis/}).

\subsubsection{Drug sensitivity analysis}\label{subsubsec_drug_sensitivity}

Drug sensitivity was predicted using the oncoPredict R package \cite{bib46}. Gene expression profiles and corresponding drug response data obtained from \url{https://osf.io/c6tfx/} were used as the training set, and bulk transcriptomic cohorts were used as the test sets. The half-maximal inhibitory concentration (IC50) of each drug was estimated for each patient sample. Associations between the risk score and predicted drug sensitivity were then evaluated using Spearman correlation analysis.

\section{Discussion}\label{sec12}

The substantial intra- and intertumoral heterogeneity of malignant cells remains a major challenge for understanding LUAD biology and improving clinical stratification. In this study, we identified nine consensus MP-defined tumor cell states shared across LUAD patients. These states exhibited distinct functional features, TF regulatory programs, microenvironmental communication patterns, progression-associated dynamics, and clinical relevance. Thus, our study provides a conserved tumor cell-state framework for interpreting malignant-cell heterogeneity in LUAD and for linking tumor-intrinsic programs to microenvironmental and clinical phenotypes.

An important finding of this study is that MP-defined tumor cell states were associated with LUAD progression and organized along divergent transcriptional trajectories. The TNF-\(\alpha\)--NF-\(\kappa\)B, MHC class II, Translation, and ATP synthesis states were more closely associated with early or less aggressive tumor cell programs, whereas Cell-cycle, EMT, and Glycolysis states were enriched during LUAD progression. Trajectory analysis further suggested that these early-stage-associated states may branch toward either proliferative or invasive/metabolic tumor cell states. These transitions were accompanied by distinct TF-regulatory dynamics. Among these TFs, FOSL2, FOXM1, and VPS72 were of particular interest. XBP1 is known to regulate differentiation and secretory programs \cite{bib47}. In particular, FOSL2 has been implicated in poor prognosis and tumor progression, partly through promoting M2 macrophage polarization \cite{bib48} and driving LUAD progression through transcriptional activation of L1CAM and the PI3K/AKT/mTOR pathway \cite{bib49}. FOXM1 is a well-known regulator of cell proliferation \cite{bib50}. Although VPS72, also known as YL1, has been reported to promote lipid synthesis in hepatocellular carcinoma by activating mTORC1 signaling \cite{bib51}, its role in LUAD remains unclear. In our data, VPS72 was enriched in malignant cells, positively correlated with UBE2C, and associated with poor prognosis, suggesting that VPS72 may represent a potential regulator of proliferative tumor states.

Our results also link MP-defined tumor cell states to distinct tumor microenvironmental contexts. Cell-cycle, EMT, and Glycolysis states showed stronger outgoing signaling toward tumor-associated macrophages, fibroblasts, and endothelial cells, suggesting that aggressive tumor cell states may more actively engage surrounding stromal and immune cells. Among the identified ligand--receptor interactions, MDK--NCL represented a broadly activated unfavorable communication axis associated with poor prognosis. This observation is consistent with previous studies showing that MDK is associated with malignant epithelial phenotypes and that MDK--NCL signaling can promote endothelial cell migration and cancer-associated fibroblast proliferation \cite{bib17,bib18,bib19}. In contrast, SCGB3A2--MARCO was more closely related to early-stage or favorable tumor states. Given that SCGB3A2 is mainly expressed in lung epithelial/Club cell programs and MARCO is involved in macrophage-associated inflammatory or phagocytic responses \cite{bib20,bib21}, this axis may reflect a distinct epithelial--macrophage communication pattern in less aggressive tumor contexts. In addition, state-specific interactions such as ANGPTL4--CDH5/CDH11 and PLAU--PLAUR were mainly associated with aggressive MP states. Previous studies have implicated ANGPTL4 in LUAD progression and EGFR-TKI resistance \cite{bib52}, supporting the potential relevance of ANGPTL4-mediated stromal interactions. These findings suggest that conserved tumor cell states are not only defined by intrinsic transcriptional programs but also linked to distinct microenvironmental niches.

The clinical relevance of these MP-defined tumor cell states was further supported by prognosis-associated cell analysis and TF-based risk modeling. Poor-prognosis-associated tumor cells were enriched in Cell-cycle, Cell adhesion, EMT, and Glycolysis states, whereas favorable-prognosis-associated cells were enriched in TNF-\(\alpha\)--NF-\(\kappa\)B and MHC class II states. By leveraging MP-state-specific TFs, we constructed an eight-TF prognostic signature that consistently stratified patient survival across multiple independent cohorts. The risk score was also associated with immune microenvironmental features and predicted drug sensitivity, suggesting that MP-state-specific TF programs may provide useful information for prognostic stratification. However, because the drug sensitivity results were based on computational prediction, they should be interpreted cautiously and require experimental validation before clinical application.

Nevertheless, several limitations should be noted. First, our analyses were mainly based on cross-sectional transcriptomic data, and dynamic state transitions were inferred computationally rather than directly observed. Second, although key TFs and ligand--receptor interactions were associated with state transitions and progression, their underlying mechanisms require experimental validation. Third, the predicted drug sensitivity analysis was based on computational inference and should be further tested in cellular, organoid, or patient-derived models. Finally, this study focused on conserved tumor cell states across LUAD patients, whereas patient-specific programs and treatment-related state changes were not systematically investigated. Despite these limitations, our study delineates conserved malignant cell states in LUAD and links them to tumor progression, microenvironmental crosstalk, TF regulation, and clinical outcome.

\section{Conclusions}\label{sec13}

This study defined nine conserved MP-dominated tumor cell states in LUAD and systematically linked these states to tumor-intrinsic programs, microenvironmental crosstalk, disease progression, TF-regulatory features, and patient prognosis. By integrating single-cell, bulk, and spatial transcriptomic analyses, we provide a tumor cell-state framework for interpreting LUAD malignant-cell heterogeneity and identifying potential biomarkers for prognostic stratification. Further experimental validation will be required to confirm the inferred regulatory mechanisms and translational relevance of these MP-defined tumor cell states.

\backmatter

\noindent\textbf{Supplementary information}

Supplementary information is available for this article, including Supplementary Figures S1--S12 and Supplementary Tables S1--S3.

\noindent\textbf{Acknowledgements}

Not applicable.

\section*{Declarations}

\begin{itemize}
\item Funding

This work was supported by the National Natural Science Foundation of China (32170675) and the Heilongjiang Provincial Postdoctoral Science Foundation (LBH-Q20147).
\item Conflict of interest/Competing interests (check journal-specific guidelines for which heading to use)

All authors declare no financial or non-financial competing interests. 
\item Ethics approval and consent to participate

No ethics approval is required.
\item Consent for publication

Not applicable.
\item Data availability 

The transcriptomic data used in this study are publicly available. The single-cell RNA sequencing data can be accessed from the Curated Cancer Cell Atlas (3CA, \url{https://www.weizmann.ac.il/sites/3CA}) and Gene Expression Omnibus (GEO, GSE189357).  The bulk RNA-seq and clinical data are available from the TCGA database and GEO repository. All datasets support the findings of this study and are openly accessible without restrictions.
\item Materials availability

Not applicable.
\item Code availability 

Not applicable.

\item Author contribution

YP designed and supervised the research. KS, YZ and RD contributed to data processing, program implementation, and manuscript preparation. YJ assisted with data collection and organized the figures and tables. All authors provided essential feedback on the final manuscript.
\end{itemize}

%%===========================================================================================%%
%% If you are submitting to one of the Nature Portfolio journals, using the eJP submission   %%
%% system, please include the references within the manuscript file itself. You may do this  %%
%% by copying the reference list from your .bbl file, paste it into the main manuscript .tex %%
%% file, and delete the associated \verb+\bibliography+ commands.                            %%
%%===========================================================================================%%

\bibliography{sn-bibliography}% common bib file
%% if required, the content of .bbl file can be included here once bbl is generated
%%\input sn-article.bbl

\end{document}

